\documentclass[11pt,letterpaper]{article}
\usepackage[T1]{fontenc}
\usepackage{lmodern,microtype}
\usepackage[margin=1in]{geometry}
\usepackage{amsmath,amssymb,amsthm,mathtools}
\usepackage{booktabs,array,longtable}
\usepackage{enumitem,xcolor}
\usepackage{listings}
\usepackage{tocloft}
\usepackage{xurl}
\usepackage[colorlinks=true,linkcolor=blue!45!black,urlcolor=blue!45!black,citecolor=blue!45!black]{hyperref}
\hypersetup{pdftitle={K215: Checked Affine Cores and a Completed CM215 Endpoint},
pdfauthor={Research packet for Vladimir Reshetnikov's ProveIt},
pdfsubject={Proof simplification, exact CM215 verification, and state-reduction experiments}}
\setlength{\cftbeforetoctitleskip}{12pt}
\setlength{\cftaftertoctitleskip}{8pt}
\setlength{\cftbeforesecskip}{3pt}
\renewcommand{\cftsecfont}{\small}
\renewcommand{\cftsecpagefont}{\small}
\setlist{itemsep=3pt,topsep=5pt}
\lstset{basicstyle=\ttfamily\footnotesize,breaklines=true,columns=fullflexible,keepspaces=true,frame=single,rulecolor=\color{black!20}}
\newtheorem{theorem}{Theorem}[section]
\newtheorem{lemma}[theorem]{Lemma}
\newtheorem{proposition}[theorem]{Proposition}
\theoremstyle{definition}\newtheorem{definition}[theorem]{Definition}
\newcolumntype{P}[1]{>{\raggedright\arraybackslash}p{#1}}
\newcommand{\N}{\mathbb N}
\newcommand{\Halts}{\operatorname{Halts}}
\newcommand{\Con}{\operatorname{Con}}
\newcommand{\Core}{\operatorname{Core}}
\newcommand{\pin}{\texttt{56e487ce8fa4cf63ad169ffe0cd2944e9fc62dd3}}
\title{K215: Checked Affine Cores\newline and a Completed CM215 Endpoint}
\author{Research and verification packet for Vladimir Reshetnikov's ProveIt}
\date{21 October 2026}
\begin{document}
\maketitle
\begin{abstract}
We study proof simplification and further state reduction for the repository's
215-state binary Turing machines for ZFC consistency. The principal new proof
artifact closes a gap in the existing affine-core proposal: it proves, in Coq,
that both literal register programs \texttt{p215} and \texttt{dagPerm} halt
exactly when the same 367-vertex accelerated graph halts, for every natural-valued
register file. The construction recognizes 1,059 drains in the 1,426-vertex
checkpoint graph, proves their exact arithmetic behavior once, checks all
concrete source bindings, and preserves halting in both directions. The
theorems require no added axioms. This is an auxiliary proof reduction, not a
reduction of the binary machine's 215 working states.

The packet also completes verification work for the existing CM215 transport
proposal, checks an inherited Lean proof-shortening patch, and supplies a
new search using a postcondition that permits spare tape data to change.
No verified machine with 214 or fewer states is claimed. The search records
distinguish bounded screening, explicit counterexamples to primitive contracts,
and actual proof-assistant theorems. All prior machine designs, proposals and
calibration results retain their attribution.
\end{abstract}

\setcounter{tocdepth}{1}
\tableofcontents
\newpage

\section{Findings, baseline and scope}

The repository baseline for this packet is \pin. The live repository changed
while the review began; this fixed revision, together with recorded Git blob
IDs, SHA-256 hashes and explicitly recorded proof patches, defines what was
actually checked. We reviewed both
the maintained proof lanes and the newly placed proposal packets. The incoming
quotient-detachment packet was also read before choosing the new work.
The primary sources are listed in Section~\ref{sec:references}.

Three distinctions are essential when evaluating progress in this program.
First, a shorter proof \emph{interface} need not be a shorter complete build.
Second, a smaller register-program graph need not compile to fewer binary
Turing-machine states. Third, a candidate that fails a uniform primitive
simulation contract has not thereby been shown to have the wrong blank-tape
consistency behavior: the failing frame might not occur in that particular
blank run. The packet uses these distinctions throughout.

\subsection{The existing exact endpoint}

The primary K215 machine has 215 working states, labels $0,\ldots,214$, a
separate halt label 215, and 430 deterministic transition rows. It uses one
two-way infinite binary tape, blank symbol 0, initial head position 0, and only
left/right moves. Its established Lean endpoint is
\begin{lstlisting}
KEndpoint215.blank_halts_iff_zfc_refutation :
  BlankHalts machine215 <-> BProv ZFCax_s [] Form.fBot
\end{lstlisting}
The companion theorem negates both sides to obtain syntactic consistency.
The Rocq lane has the analogous exact dense-machine endpoint. These are
external checked equivalences with finite ZFC refutability. They do not prove
unconditional nonhalting or the consistency of ZFC, and they are not a new
construction of an internally encoded PA or ZFC proof object. The ordinary
argument that the constructions and equivalence can be formalized in ZFC is
inherited from the project's calibration work.

\subsection{What was already available}

The current proposals already contain several improvements that should not be
advertised again as new results: a single checked certificate plus a generic
mirror theorem; quotienting addresses by their trailing zero suffix; reducing
2,240 checkpoint entries to 1,426; balanced certificate lookup; automatic
certificate inference from literal trees; shared decoder contexts; concrete
zero-credit parser and prefix-credit theorems; and local 215-state transport
proofs. The quotient-detachment packet already proves a larger numerical
inference rule equiconsistent with ZFC, but reports no smaller physical
machine. The existing affine-core packet already identified the 1,059 drains,
the 367-vertex graph and 33 shared update matrices. Its generic drain theorem,
source-to-graph binding and complete contraction proof were not kernel checked
there. Those are the main gaps closed here.

The maintained Lean source proof already uses a balanced \texttt{CTree}
certificate lookup and staged concrete kernel checks. The earlier balanced
lookup proposal addresses the Rocq flat-list checker; it should not be described
as a missing Lean optimization. The existing Lean automatic-certificate
proposal is a different improvement: infer the relation from the literal trees,
then check a normalized 1,426-entry certificate with shared tree bindings.
Attaching that existing facade to the complete endpoint is a useful separate
integration task. Neither that proposal nor the timings in this packet establish
which operation dominates a full Lean build.

\subsection{Verification record}

The accompanying component receipts contain exact commands, compiler versions,
source hashes, module outcomes and assumption reports. The following inserted
record is generated from the final accepted runs, rather than from a planned
test count.

\begin{itemize}
\item Coq 8.18.0 compiles all seventeen source-lane modules, including the
exact two literal source/core equivalences, same-name public facade and typed
client. All twenty-five printed assumption reports are closed. A separate
recursive \texttt{coqchk EndpointClient} passes with exit code zero.
Compilation takes 92.585 seconds and the recursive check 440.022 seconds in
the recorded run; these are reproduction observations, not comparative
speed claims.
\item Lean 4.32.0 accepts the complete 184-module K215 project closure and
both original typed endpoint audits. This baseline contains 180 byte-identical
project modules and four disclosed certificate-proof batching overlays.
A subsequent verified resume preserves the source, inventory, toolchain and
proof receipts and performs zero new compiler attempts.
\item The exact CM215 attachment passes all eight modules: seven inherited
modules and an independent typed client. The client checks finite-run halting
and its negation against the concrete machine. All printed dependency axioms
are among \texttt{propext}, \texttt{Classical.choice} and \texttt{Quot.sound}.
\item The unchanged inherited \texttt{KSeedModel222} patch passes all ten
affected project modules and four endpoint/audit clients, including K215 and
CM215 endpoints. With this patch and the batching overlays together, 179 of
the 184 project modules remain byte-identical to the pin and five have
disclosed proof changes. Printed axioms remain within the same standard set.
\item Independent finite audits replay all 59 survivors of the relaxed-spare
screen, all 639 coordinated-provider cases and parent controls, and the
three detailed startup/prefix checks. Their distinct failure and timeout
categories are reported in Section~\ref{sec:214}. The separate symbolic verifier accepts
all 1,900 increment occurrences for the rejected no-append table; this is
quantified symbolic evidence outside a proof-assistant kernel.
\end{itemize}


\section{A proof architecture that separates the right obligations}

The most useful long-term structure has three semantic interfaces. The first
relates finite derivations in the numerical calculus to finite ZFC
refutations. The second relates the literal register program to a compact
graph of tested operations and total arithmetic macros. The third relates
that program, or a previously proved exact machine, to the literal binary
transition table. Each interface has different optimization opportunities.

At the logical interface, use one interpretation/substitution theorem and
derive special cases by composition. At the register-program interface,
compute finite correspondence data from literal syntax and validate it in the
kernel. At the machine interface, use local invariants and positive finite
simulation blocks when the new table is a small modification of a proved table.
Rebuilding the calculus or the whole register runtime for such a local table
change adds avoidable work.

A concrete dependency cleanup is available in the source imports.
\texttt{Prog215.v} uses \texttt{K230.Prog230} only for its small generic
\texttt{cedge} datatype, constructors and \texttt{cv} decoder; it does not
reference \texttt{p230}. That inherited module contains 5,166 lines and took
19.1 seconds in the focused build recorded here. Extracting the generic codec
into a shared module, updating generated imports and preserving required
public names would remove the unrelated literal 230-state program from this
source closure. This is a concrete follow-up proposal, not an applied patch
or a measured saving for the full endpoint. The current proof deliberately
retains byte-identical production program inputs.

The Lean build also motivates a resource-oriented proof change. The original
monolithic \texttt{fwd\_ok12} check terminated with exit code 137 under
an 8\,GiB memory limit, without a timeout. The recorded overlay replaces the proofs in certificate
chunks 12--15: each 140-entry check becomes four private 35-entry
\texttt{decide +kernel} lemmas, followed by checked Boolean reassembly.
Imports, public statements, program data and certificate data are unchanged.
The companion generator patch emits exactly the four checked replacements
and leaves the other fourteen partitioned source modules byte-identical.
This applies the quarter-block technique already present in the earlier
normalized and automatic-certificate proposals. The source exceptions and
the failed original attempt are preserved explicitly. This adjustment makes
the checking work smaller per lemma; the four modules grow from 11 to 36 lines
each, adding 100 generated lines, and
must not be counted as part of the separate handwritten line reduction in
\texttt{KSeedModel222}.

For the two current register programs, normalized checkpoint correspondence
already gives a particularly direct proof. Our affine construction should not
automatically replace that fast path. Its main benefit is a reusable semantic
interface for future sources with different implementations of transfer, copy
or clearing loops. Such programs can have unequal checkpoint graphs while
performing the same total affine operations. A proof about their arithmetic
cores can then avoid repeating hundreds of low-level loop inductions.

\subsection{The next substantial runtime-proof reduction}

The most promising next reduction in the large runtime proof is to complete
the attachment of the existing framed decoder component. That earlier packet
already supplies an arbitrary-tape framing theorem and checks 281 local blocks
over 234 decoder roles, with 97 exposed primitive-entry contexts. These are
shared internal traversals of the source DAG. They are not yet 97 complete
primitive contracts or a replacement for the 4,751-occurrence site-coverage
proof. The missing attachment has three concrete parts.

\begin{enumerate}
\item Bind the decoder's source DAG expansion to the literal
\texttt{KSource215} program inside the proof assistant. The existing pilot
has an independent Python expansion audit, which is useful evidence but does
not supply this kernel equality.
\item Prove by induction on a valid source path that decoder roles produce
the framed block-chain premises. Carry the exact head location and preceding
bits, including the single left marker at the root. Then invoke the existing
arbitrary-tape theorem, so the unexamined gap, header, registers and spare
suffix are preserved without repeating whole-tape calculations at each site.
\item Prove the exposed primitive and return contracts with a variable code
prefix, the exact prefix-language restriction, arbitrary register data and
the appropriate zero/nonzero guard. Factor carry and strip behavior into
list-induction lemmas, compose the existing register-box contracts, and
recover the same public site-coverage theorem.
\end{enumerate}

The prefix-language condition is essential. Two preceding bits suffice for
the already checked internal decoder blocks; they have not been proved
sufficient for arbitrary leaf actions and pointer returns. Break and
zero-decrement returns can overflow the code prefix at particular source
sites, so the marker movement and ancestor counts must remain in the contract.
Likewise, a shared physical state alone is not a sound decoder role: the source
DAG node is part of its meaning. These constraints explain precisely what
must be proved before the existing local pilot can simplify the full runtime.

This is a proposed integration of prior work, not a completed replacement in
this packet. It has greater potential to remove repeated runtime proof work
than merely lowering source-tree depth, but its effect on kernel memory and
clean-build time must be measured after the exact endpoint is reattached.
The current packet retains and checks the full site-coverage dependency chain.

\section{The exact checkpoint semantics}

A register file is a function $R:\N\to\N$. Although the concrete program uses
nine registers, the generic theorems quantify over every natural-valued register
file. Unused registers are preserved. Pointwise equality
\[
 R\equiv S\quad\Longleftrightarrow\quad (\forall i\in\N)\ R_i=S_i
\]
is carried explicitly, so functional extensionality is not an added axiom.

Each graph vertex is either a halt or a test
\[
 v:\operatorname{Test}(r,(u,a),(w,b)),
\]
where $a$ and $b$ are finitely supported nonnegative increment vectors. If
$R_r>0$, the graph first decrements $r$, then adds $a$, and continues at $u$.
If $R_r=0$, it adds $b$ and continues at $w$. The order of the decrement and
the following arithmetic update is part of the semantics. Omitting that
decrement would invalidate the contraction even when every target vertex was
correct.

The source-to-graph checker verifies the literal instruction at each graph
address, the register being tested, both branch directions, every increment
count between tests, and the next source address modulo trailing zeros.
It uses the inherited segment executor, which is already proved to execute
only increments and control instructions until the next test or halt. A
finite graph-closure check guarantees that each target is defined.

The checker does not trust the Python graph generator, an address hash, or a
claim that two trees look alike. Its concrete acceptance theorems are evaluated
against the literal \texttt{Prog215.p215} and \texttt{DagPerm.dagPerm}
definitions. Their bytes match the maintained production files at the pinned
revision. The generated data is input to a proof, not an extra axiom.

\section{The total drain theorem}

\begin{definition}
A drain at $v$ has the shape
\[
 v:\operatorname{Test}(r,(v,a),(w,b)),\qquad a_r=0.
\]
Thus the positive branch returns to $v$ without replenishing the tested
register. The zero branch supplies the exit and its fixed increments.
\end{definition}

\begin{theorem}[Exact all-register drain semantics]\label{thm:drain}
For every register file $R$, a drain at $v$ reaches $w$ in exactly $R_r+1$
checkpoint transitions. Its final register file $D(R)$ is
\[
 D(R)_i=
 \begin{cases}
 b_r,&i=r,\\
 R_i+R_r a_i+b_i,&i\ne r.
 \end{cases}
\]
In particular, the exit increments are applied exactly once, including a
possibly nonzero increment of the tested register on exit.
\end{theorem}

\begin{proof}
Put $n=R_r$ and induct on $n$. For $n=0$, the zero edge is taken immediately,
giving the formula and one transition. For $n+1$, the positive edge changes
the tested register to $n$, because $a_r=0$, and adds $a_i$ to every other
register. Applying the induction hypothesis to this new register file yields
$R_i+a_i+n a_i+b_i=R_i+(n+1)a_i+b_i$ for $i\ne r$. The tested register has
value $b_r$ at the exit. The number of transitions is $1+(n+1)=n+2$.
\end{proof}

The checked theorem is \texttt{DrainCore.drain\_exact\_steps}. The supporting
\texttt{drain\_halts\_iff} theorem proves equivalence of halting at the original
configuration and the accelerated exit configuration. No bound on $R_r$ is
used. The proof also permits repeated entries in the finite increment list:
their contributions are summed by the definition of \texttt{amount}.

\subsection{Why this is a useful simplification}

The loop proof depends on a tiny syntactic condition and one natural-number
induction. It does not depend on where the loop occurs, on the numerical
calculus, on the caller's tape encoding, or on a prior bound on the registers.
All 1,059 occurrences can therefore use the same theorem. Their certificates
check the loop shape and exit structure; they do not contain 1,059 separate
arithmetical arguments.

The earlier proposal encodes these operations as homogeneous $10\times10$
matrices. The present kernel theorem uses the direct arithmetic formula and
composition by \texttt{normalize}. Its trusted proof does not require the
precomputed dictionary of 33 matrices. The matrix presentation remains a useful
human-readable and synthesis representation, but the matrix dictionary's
existing Python validation should not be confused with our new kernel theorem.

\section{Contracting drains without losing divergence}

Let $\operatorname{norm}_f(v,R)$ follow at most $f$ recognized drains,
accelerating each by Theorem~\ref{thm:drain}. If the current vertex is not a
drain, normalization stops immediately. Define the accelerated transition from
any nonhalting vertex by executing \emph{one genuine original tested edge} and
then normalizing its successor. This first real transition guarantees positive
progress even when there are no following drains.

\begin{theorem}[Generic contraction correctness]\label{thm:core}
For every graph $G$, every fuel $f$, every vertex $v$ and every register file
$R$,
\[
 \Halts_{\Core_f(G)}(v,R)\quad\Longleftrightarrow\quad\Halts_G(v,R).
\]
The analogous statement holds after normalizing the initial configuration.
\end{theorem}

\begin{proof}
For soundness, expand each accelerated edge into its original first transition
and the finite drain runs. Normalization preserves halting by induction on its
fuel. Composing the finitely many expanded edges gives an original halting run.

For completeness, use an original halting-time budget. Define
$\operatorname{Within}_G(n,v,R)$ to mean that a halt can be reached within $n$
transitions, allowing an already halted configuration at any budget. An original
transition consumes at least one unit of a nonzero budget. Determinism, up to
pointwise register equality, implies that following an actual prefix cannot
increase the remaining budget. Each drain is such a finite prefix. Consequently,
normalization cannot increase the budget. After one genuine tested transition
and normalization, induction applies to a strictly smaller budget. This yields
a finite accelerated halting derivation.
\end{proof}

The generic theorem is valid for any normalization fuel. A rank bound is needed
to prove that the accelerated control flow stays in a chosen smaller core, not
to establish the generic halting equivalence. This separation makes the proof
more robust: an insufficient fuel value cannot silently certify an incorrect
halting equivalence, although it may fail to establish the advertised core size.

For the concrete graph, every removed vertex has a strictly decreasing exit
rank, with maximum rank six. Kernel computation checks that normalizing every
defined vertex for six passes reaches the retained set. It also checks closure
of every accelerated successor in that set. The retained vertices are distinct;
there are 366 tests and one halt. The initial vertex 0 is already retained.

\section{Binding the 367-vertex core to both literal programs}

\begin{theorem}[Literal source adequacy]\label{thm:source}
For every natural-valued register file $R$,
\begin{align*}
 \operatorname{rm\_halts}(\texttt{p215},([],R))
 &\Longleftrightarrow \operatorname{core\_halts}(6,G,0,R),\\
 \operatorname{rm\_halts}(\texttt{dagPerm},([],R))
 &\Longleftrightarrow \operatorname{core\_halts}(6,G,0,R),
\end{align*}
where $G$ is the literal checked checkpoint graph supplied in the packet.
\end{theorem}

\begin{proof}
The finite \texttt{node\_ok} checks establish source instruction, branch,
increment and address agreement at all 1,426 graph vertices for each literal
program. Graph closure supplies all successor vertices. A generic adequacy
proof converts a graph halting derivation into source segments, and conversely
uses strong induction on the source halting time to peel off a tested
instruction and its following increment-only segment. Address equivalence
handles the implicit infinite zero suffix. Pointwise register equalities account
for the effects of the segment. The checked initial addresses are equivalent
to the empty path. Apply Theorem~\ref{thm:core}.
\end{proof}

\begin{samepage}
The exported concrete theorems are:
\begin{lstlisting}
p215_core_halts R :
  rm_halts p215 ([], R) <-> core_halts 6 k215 0%N R

dagPerm_core_halts R :
  rm_halts dagPerm ([], R) <-> core_halts 6 k215 0%N R

p215_dagPerm_affine : forall R,
  rm_halts p215 ([], R) <-> rm_halts dagPerm ([], R)
\end{lstlisting}
\end{samepage}
The last theorem has the existing source-correspondence shape. It is derived
by rewriting both sides through the common core. The five new modules separate
generic drain semantics, generic contraction, finite graph data, generic source
adequacy, and the two concrete source bindings. The final source proof includes
ten inherited prerequisite modules. No complete production Rocq binary/ZFC
rebuild is inferred merely from this new source theorem.

\begin{center}
\begin{tabular}{lr}
\toprule
Quantity & Checked value\\\midrule
Original checkpoint vertices & 1,426\\
Recognized drain vertices & 1,059\\
Retained core vertices & 367\\
Retained tests / halts & 366 / 1\\
Retained outgoing edges & 732\\
Maximum normalization rank & 6\\
Local drain shapes, prior analysis & 24\\
Affine matrices, prior presentation & 33\\
\bottomrule
\end{tabular}
\end{center}

The retained control interface is approximately $25.74\%$ of the original
checkpoint set. This number is not a measured compilation-time improvement.
The accelerated relation still consults the original graph data; the
367-control claim is established by the retained-set and closure checks.
All removed shapes must still be checked, the graph is explicitly bound to two
programs, and the generic library has a cost. For today's already equal source
graphs, the earlier direct normalized certificate may remain faster and smaller.
In this implementation, the generated graph/rank and location sources total
about 285\,KB, compared with 109,958 bytes for the earlier normalized
certificate. Thus the core reduction is not a generated-byte reduction.
The new result becomes more valuable when comparing semantically equivalent
sources with different low-level loop decompositions.

\section{Which other 215-state machine is easiest to prove?}

The answer depends on the proof obligation being minimized. For the least
additional work from the current certified K215 endpoint, CM215 is the strongest
practical choice among the inspected proposals. For a smaller isolated marker
argument, M124/FM215 is competitive. For clearer independent scheduler
mathematics, the zero-credit source family is attractive, but its exact
binary-machine endpoint still needs substantially more work.

\begin{center}\small
\begin{tabular}{P{0.16\textwidth}P{0.35\textwidth}P{0.39\textwidth}}
\toprule
Candidate & Main proof advantage & Remaining cost or qualification\\\midrule
CM215 & Same source; six-row local transport with an existing exact dense adapter & Reuses the entire K215 calibration; smallest incremental endpoint route, not a smaller total logical foundation\\
M124 / FM215 & Four changed rows and two local marker facts & Existing standalone transport still needs an exact repository-semantics adapter\\
ZC215 & Occurrence-injection proof removes protected-tail credit & Concrete scheduler and exact physical implementation must still be connected end to end\\
R41 / rebalanced & Decoder depth decreases from 47 to 41 & All 4,751 primitive occurrences and 6,176 guarded cases remain; no measured full-proof reduction\\
Context215 & Five changed rows & Six block simulations and history conditions; row count understates proof cost\\
DR215 / GC-ZC215 & More substantial source design changes & New parameter/selector or generalization/scheduling obligations make relative transport longer\\
\bottomrule
\end{tabular}\end{center}

\begin{samepage}
These comparisons concern the currently available proof routes, not lower
bounds on proof complexity. All listed candidates were already present in
the reviewed proposals.

\subsection{The exact CM215 patch}

In dense labels, the existing CM215 table changes exactly the following rows.
Every unlisted row is inherited from the primary K215 table.
\begin{center}
\begin{tabular}{ccc}
\toprule
State / read & K215 & CM215\\\midrule
52 / 0 & $0\,R\,149$ & $0\,L\,170$\\
151 / 1 & $1\,R\,170$ & $1\,R\,52$\\
160 / 0 & $0\,R\,170$ & $0\,R\,52$\\
170 / 0 & $0\,R\,163$ & $0\,R\,178$\\
170 / 1 & $1\,R\,145$ & $1\,R\,125$\\
184 / 0 & $0\,R\,170$ & $0\,R\,52$\\
\bottomrule
\end{tabular}
\end{center}
\end{samepage}
The old-run invariant has three clauses, with $h$ the head position and $t$ the
tape:
\[
 q=52\Rightarrow t(h-1)=1,\qquad
 q=151\Rightarrow t(h+1)=1,\qquad
 (q=170\land t(h)=0)\Rightarrow t(h-1)=0.
\]
The boundary map sends old state 170 to new state 52 and fixes every other
label. Most old steps correspond to one new step. The exceptional blocks are
$52\to170\to125\to149$ and $52\to170\to178\to163$, each of length three.
Every simulated boundary preserves the whole tape and head position.
New state 170 is still used inside these blocks, so identifying two old
boundary controls does not remove a physical working state.

Because the boundary map is noninjective, its inverse cannot simply be used to
prove the reverse implication. The inherited positive-stuttering theorem uses
determinism and positive block lengths to synchronize any finite target halt
with the source run. This is the mathematical reason the small transport proof
can replace a new register-runtime and startup verification.

The dense CM215 table has SHA-256
\begin{quote}\ttfamily\small
\nolinkurl{1f4f084f7293e2b042d5746039f793ff72ed429083b504970c939a5515e666cf}.
\end{quote}
The parent dense table has SHA-256
\begin{quote}\ttfamily\small
\nolinkurl{fdbcef0825886035c3091c8537a69535bc96f9e265f0ae45863543d40ae70090}.
\end{quote}
The exact 430 rows of both tables are independently compared with the Lean
definitions. The new contribution is the completion and audit of the concrete
endpoint attachment against the pinned current dependency closure; the table,
invariant, transport theorem and attachment design come from the earlier
\texttt{k215-transport} proposal.

\section{A direct opportunity to shorten the Lean semantic proof}

The earlier affine-core packet proposes a shorter proof in
\texttt{KSeedModel222}. It removes three repeated formula inductions by using an
already proved semantic encode/decode theorem at the model boundary. A proof that descends
through every formula constructor should be written at the interpretation
boundary; callers should not repeat it when they only need a special case.

The concrete change removes \texttt{enc\_pos}, \texttt{enc\_dec\_enc} and
\texttt{eqFree\_dec\_enc}. These were used to take a named formula through
encoding, decoding, equality-freeness and re-encoding before applying the
bundle rule. Instead, construct the equality-free named bundle $\varphi$
directly and prove that its interpreted code is derivable using
\texttt{C222.bundle}. The existing hypothesis about a model of all derivable
named formulas then gives $\operatorname{Sat}(\varphi)$. Apply the already proved
semantic round trip
\[
 \operatorname{Sat}_{\mathrm{mem},e}
       (\operatorname{dec}(\operatorname{enc}(\varphi)))
 \ \Longleftrightarrow\
 \operatorname{Sat}_{\mathrm{mem},e}(\varphi)
\]
from \texttt{Code254Models.sat\_encode\_decode}. Its existing irreflexivity
hypothesis handles the bottom-formula alias. No new interpretation theorem or
import is needed. The one remaining numerical lower bound follows from
\texttt{le\_P\_right} and $F=1$, rather than induction on a formula.

On the pinned source, the module decreases from 269 to 234 lines
(35 lines, $13.0\%$) and from 11,711 to 10,436 bytes
(1,275 bytes, $10.9\%$). The patch changes no machine table or source program.
The public \texttt{seed\_eqFree} theorem and the endpoint statements retain
their types. The helper API does change: \texttt{derivable\_bundle} becomes
\texttt{derivable\_named\_bundle}, and the three redundant helpers are removed.
Any additional consumers outside the checked endpoint closure would need to
use the new interface or retain a compatibility layer.

The verification component rebuilds the exact ten affected project modules
and four endpoint clients in a separate overlay after checking the declared
184-module baseline, including the four proof-only batching overlays.
This makes the acceptance criterion dependent theorem compilation,
not merely a text diff or an isolated tactic check. The final outcomes and
assumption reports are recorded in the accompanying receipts.


This is a different kind of improvement from the graph contraction. It can
reduce maintained handwritten source directly, while the graph construction
adds generic infrastructure and checked generated data. Both are useful, but
their measurements should not be mixed. In particular, lines removed from one
module do not establish a proportional improvement in total endpoint build
time.

\section{Research toward 214 states}\label{sec:214}

\subsection{What was searched}

We explored new local quotient shapes in addition to the existing same-head,
same-tape folds: identifying states with a head displacement, and identifying
them while toggling a nearby marker cell. The scripts record their finite
window, step bounds and chosen invariant. Their failures only delimit these
specific proof/search families.

The main new experiment addresses the most promising relaxation suggested by
the automatic-certificates packet. The register runtime's logical interface
uses nine registers, a continuation, a legal gap, and a nonempty spare suffix.
The high-level representation need not expose the exact spare suffix returned
by each primitive. Requiring an exact whole-tape join can therefore reject
transformations that preserve every logical register and continuation.

Our relaxed decoder accepts any valid nonempty spare suffix and legal gap,
while requiring the nine logical registers and continuation to be exact. The
search deletes one state, rewrites its unique external incoming row, and tests
the new machine on valid primitive frames. For the main survey, 29 covered
labels, 214 replacement targets, two writes and two directions give
\[
 29\cdot214\cdot2\cdot2=24,824
\]
candidates. The initial screening cap is 5,000 transitions. It is a bounded
screen, not a decision procedure for arbitrary candidate correctness.

\subsection{Counterexamples for all short-screen survivors}

Fifty-nine candidates survive the first short opcode-class tests. They are
then tested at actual complete-source primitive occurrences. Fifty have an
explicit wrong decoded logical result at a clean dispatcher. The other nine
initially exceeded a 20,000-transition checking cap. A separate cycle analysis
then found literal repeated complete configurations, with zero head
displacement, by steps 585--609. No halt or successful required redispatch lies
between each recorded repeat. Determinism therefore proves infinite repetition
on those particular valid primitive inputs.

An independent byte-array simulator checks all 59 witnesses, the corresponding
parent frames, the expected source effects, and the claimed first redispatch
times. It does not import the search implementation. For each cycle it compares
the state, head and every nonzero tape cell. The finite witness plus
determinism establishes the claimed loop without extrapolating from a timeout.

These are counterexamples to the uniform relaxed primitive contract that
requires the correct result at the first clean redispatch. They do not exclude
a different simulation relation that permits extra intervening redispatches. They do
not prove that every mutated blank-start machine has a different consistency
behavior. Nor do they turn every one of the earlier bounded-screen failures
into a complete mathematical refutation. The justified conclusion is that this
pipeline produced no successful reduction and that its 59 initially promising
survivors all fail the intended reusable simulation obligation.

\subsection{A concrete shared-helper obstruction}

Dense state 189 is reached externally from row $66/0$. In the inspected
implementation it participates both in spare extension for increments and in
restoration on a zero decrement. A replacement that looks correct when the
increment is allowed to leave a different spare suffix can still destroy the
decrement restoration route. This is why weakening the spare-output condition
alone does not automatically free a state: shared physical control has more
than one logical role.

There is a positive, quantified sub-contract. Change the native row
\[
 66/0:\quad 1\,R\,198\quad\longmapsto\quad 0\,L\,54
\]
and delete native state 198 (dense state 189). The resulting research table has
214 listed working labels and 428 rows; 212 working labels remain reachable
in its transition graph. The inherited symbolic runtime checker proves
all 1,900 literal increment occurrences with the exact logical increment and
continuation, while leaving the arbitrary nonempty spare suffix unchanged.
Its induction rules cover arbitrary natural register values, arbitrary natural
gap and arbitrary finite spare data; this is stronger than bounded testing.
The run uses 672,539 symbolic rules. It is an ordinary quantified symbolic
argument checked by a Python verifier, not a proof-assistant kernel theorem.

The same table fails a concrete zero-decrement contract. At zero-based source
index 2 (the third leaf),
address \texttt{0100}, command \texttt{CD0}, take all nine registers zero,
spare suffix $[0]$ and gap zero. The parent returns after 640 transitions with
all logical registers zero and continuation one-bits at positions $[1,2]$.
The research table first returns after 3,966 transitions with continuation
one-bits at $[1,3,4,5]$ and register values
\[
 [0,0,0,0,0,1,0,0,2].
\]
An independent simulator replays this mismatch. Thus the 214-state table is
delivered only as a rejected uniform-runtime experiment, with its failed
obligation prominently identified. It is not delivered as a verified solution
to the consistency machine problem; its blank-run calibration remains
unestablished.

All 856 direct write/direction/target replacements for the incoming row also
fail the stated relaxed short screen. A targeted 856-choice provider experiment
likewise finds no sampled full-contract survivor. These are bounded
observations, not global impossibility results for a joint rewrite of the
shared region.

\subsection{Coordinated reuse of an existing provider}

A further fixed family redirects dense row $66/0$ to another provider $p$ and
copies either or both outgoing rows of state 189 into $p$, then deletes 189.
Excluding providers 66 and 189 gives $213\cdot3=639$ tables. All 109 survivors
of the short screen fail the subsequent complete-source sample refinement.
An independent simulator audits the whole family, with the following results.

\begin{center}
\begin{tabular}{lr}
\toprule
Observed obstruction or remaining bound & Tables\\\midrule
Wrong first-redispatch output & 338\\
Outward motion into the infinite blank exterior & 223\\
Literal repeated configuration & 5\\
Unexpected halt during a nonhalting primitive & 2\\
Unresolved at the 20,000-transition bound & 71\\\bottomrule
\end{tabular}
\end{center}

Thus 568 tables have explicit primitive-contract obstructions; the remaining
71 have only bounded failures. The ray checks verify a read-zero self-loop
moving away from every nonzero cell, and the cycle checks compare complete
configurations. These conclusions are stronger than timeouts but remain
specific to the stated primitive contract.

Three variants, using providers 95, 188 and 213 with both rows replaced, are
particularly informative. Each has 214 listed and graph-reachable working
states. Their first decoded blank-start frames match the parent's complete
tape and head: at 79,633 transitions for the first two, and 79,639 for the
third. This experimental first decoded frame is distinct from the maintained
formal bootstrap endpoint. They pass respectively 5,944, 5,936 and 5,918
guarded source cases before their first failed case. None of those finite
successes establishes correctness.

Their failures occur at zero-based source indices 4,578, 4,572 and 4,560.
The bad prefixes read only positions $-10$ through 14, 15 or 21 relative to
the initial head. Every input cell depending on registers, gap or spare data
lies at position at least 64. The returned continuation differs from both
branches of the source decrement. The same finite prefix therefore fails for
every natural register file and every legal gap and spare suffix at that
address. A rescue by a stronger input invariant must exclude the address
itself; excluding only the displayed zero-register valuation is insufficient.
The independently replayed tables and witnesses are retained as rejected
experiments in \texttt{state\_search/forced\_provider189/}.

The addresses also have a useful source-level explanation. They lie in the
mandatory final gate of the AX routine: the first comparison, restoration of
its saved value, and the second comparison. Direct graph replay reaches a
\texttt{Clean} AX entry after 180 checkpoint transitions from the supplied
initial decoded frame. The inherited total macro-run theorems then reach the
gate with registers $[7,T,W,0,0,0,0,0,0]$ for finite $T,W$. The literal graph
visits source indices 4,560, 4,572 and 4,578 at offsets $1$, $T+4$ and $T+6$
from gate entry. Thus an ordinary argument from inherited totality and the
checked literal mapping establishes their eventual occurrence in the parent
source run; a long bounded trace was merely still doing earlier arithmetic.
This argument is documented in
\texttt{state\_search/source\_reachability/}. It is not a new kernel
reachability theorem, and it does not establish that any altered physical
machine follows the parent to those controls or has a different blank-tape
halting behavior.

The overwritten providers are source-address decoder controls. Provider 188
is especially localized: it separates two closed transfer loops, from
register 7 to 6 and from register 4 to 5. Three simple swaps or reassociations
of this region preserve 202 distinct source DAG nodes; the corresponding
literal two-node rewires each still have 215 working states and no ordinary
Mealy-state merge. The source transformations are justified by closed-block
sequencing and the disjoint register pairs. They do not certify those physical
rewires: a changed intermediate decoder context still needs its own runtime
proof. This small experiment saves no state. A more substantial joint
implementation of the two transfers, their decoder and the shared
append/restoration interface remains a concrete synthesis target.


Subsequent synthesis should require the increment and zero-decrement
obligations together from the beginning. Source equivalence, exact decoder
binding, quantified primitive contracts and positive simulation progress are
separate acceptance gates. Validating only an increment subcase and counting
a deleted shared state does not establish the desired calibration.

\section{Integration plan}

The following order makes each change reviewable and keeps the exact existing
endpoint available throughout.
\begin{enumerate}
\item Apply the minimal build-driver correction if using the supplied current
Lean driver. Its compiler identity query should use quiet Lake output: one-time
manifest diagnostics otherwise change the captured identity string even when
the compiler binary and all hashes are unchanged. The packet retains the
reproduced failure, the one-line patch, and the resumed-build check.
Apply the separately disclosed four certificate-proof batching patches when
using the recorded memory-limited build, together with their checked generator
postpass so regeneration preserves the proof structure.
\item Integrate the checked \texttt{KSeedModel222} source patch and its
dependent theorem checks, subject to the precise outcome in the included
verification record. The patch changes no generator or machine table.
\item Add CM215's already designed exact-table adapter and endpoint with the
completed verification receipts. Keep the six changed rows, invariant clauses,
literal-table bindings and attribution together.
\item Add the affine-core files as a separate checked source lane. Its exact
public source-equivalence theorem can be used behind a same-statement facade.
Do not remove the existing normalized certificate merely because the retained
core has fewer vertices. Measure clean compilation and review both proof sizes
before choosing a maintained default.
\item Use the core interface for new source designs, while keeping binary
allocation and quantified runtime verification as separate gates. A successful
semantic source simplification is a proposal for physical synthesis, not an
automatic state saving.
\end{enumerate}

The archive is additive and does not modify the remote repository. It preserves
the source inputs needed for offline reproduction, compiler logs and small
patches. Compiled proof objects and toolchain binaries are intentionally not
part of the deliverable; the build scripts reconstruct the relevant objects
from the supplied sources.

\section{Further research questions}

\begin{enumerate}
\item \textbf{Can different source variants share a proved arithmetic core?}
The present two programs already have the same checkpoint graph. Apply the
new interface to zero-credit, transfer, or rebalanced sources whose low-level
graphs differ. A successful common-core certificate would demonstrate the
benefit that motivates this additional infrastructure.
\item \textbf{Can more total loop patterns be recognized soundly?}
The drain rule requires no invariant beyond its syntax. Useful extensions
include paired copy-and-restore loops, swaps, and guarded affine blocks.
Every extension should carry an exact termination measure and should preserve
the explicit decrement at a retained test.
\item \textbf{Can a verified coefficient representation make the data smaller?}
The direct arithmetic normalizer avoids trusting the prior matrix dictionary.
A kernel-checked sparse matrix evaluator with a shared dictionary could reduce
duplicated arithmetic data, provided measurements justify the added layer.
\item \textbf{Can M124/FM215 receive the same exact endpoint attachment?}
Its local marker argument is smaller than CM215's. A short adapter into the
repository's dense semantics may make it the preferable next independently
attached 215-state variant.
\item \textbf{Can the increment and zero-decrement helper be resynthesized
together under existential spare output?}
The concrete failures identify a shared restoration obligation that a local
increment-only optimization misses. This joint region is a more informative
target than another broad same-tape state-fold census.
\item \textbf{Can equiconsistency simplify the source enough to survive
allocation?}
The quotient-detachment result demonstrates a mathematically valid relaxation
of the numerical calculus, but its first physical implementations cost more
states. Search objectives must combine arithmetic proof simplicity, address
sharing, command selection, spare management and physical return routing.
\item \textbf{Which simplifications reduce total verification cost?}
Measure handwritten proof lines, generated bytes, kernel memory, clean-build
time and independent checking separately. A single headline reduction in
certificate entries or graph vertices does not answer all five questions.
\end{enumerate}

\section{Reproduction and provenance}

The top-level README provides exact build commands for each component. The
manifest binds every delivered payload except the manifest itself. A manifest
check establishes file identity; it is not a semantic theorem. The default
portable verifier also replays the independent search witnesses. The Coq and
Lean build commands are separate so that a Python integrity check cannot be
mistaken for a kernel check.

The mathematical calibration, primary binary table, literal register programs,
segment semantics, normalized-address tools, CM215 design, and incoming
three-induction proposal are inherited. New contributions are the checked drain
and contraction library, the generic checked graph/source bridge and its exact
two-program instantiation, the completed verification work explicitly recorded
in this packet, the minimal driver fix, and the new local-search experiments
and independent counterexample audit. The packet does not claim a new logical
independence theorem, a new proof of ZFC consistency, a minimum-state result,
or a verified 214-state machine.

\clearpage
\section{Primary sources}\label{sec:references}

All repository links below are pinned to the reviewed revision. The original
files and packet notices identify their individual contributors and inherited
licenses; the archive retains the relevant notices. References to the new
proofs refer to source files in this archive and their final verification
receipts.

\begin{enumerate}[label={[\arabic*]}]
\item ProveIt, primary K215 description and exact machine convention.\newline
\href{https://github.com/VladimirReshetnikov/ProveIt/blob/56e487ce8fa4cf63ad169ffe0cd2944e9fc62dd3/Computability/BusyBeaver/Research/ZFC/candidate-k215/README.md}{\texttt{candidate-k215/README.md}}.
\item ProveIt, exact Lean endpoint,
\href{https://github.com/VladimirReshetnikov/ProveIt/blob/56e487ce8fa4cf63ad169ffe0cd2944e9fc62dd3/Computability/BusyBeaver/Lean/BusyBeaver/ZFC/KEndpoint215.lean}{\texttt{KEndpoint215.lean}}.
\item ProveIt, current simplification and proposal index,
\href{https://github.com/VladimirReshetnikov/ProveIt/blob/56e487ce8fa4cf63ad169ffe0cd2944e9fc62dd3/Computability/BusyBeaver/Research/ZFC/k225-formal/verification/k215/simplification/proposals/README.md}{\texttt{simplification/proposals/README.md}}.
\item ProveIt, prior affine-core proposal,
\href{https://github.com/VladimirReshetnikov/ProveIt/tree/56e487ce8fa4cf63ad169ffe0cd2944e9fc62dd3/Computability/BusyBeaver/Research/ZFC/k225-formal/verification/k215/simplification/proposals/affine-cores}{\texttt{proposals/affine-cores/}}.
\item ProveIt, existing local transport and endpoint attachment sources.\newline
\href{https://github.com/VladimirReshetnikov/ProveIt/tree/56e487ce8fa4cf63ad169ffe0cd2944e9fc62dd3/Computability/BusyBeaver/Research/ZFC/k225-formal/verification/k215/simplification/proposals/k215-transport}{\texttt{proposals/k215-transport/}}.
\item ProveIt, automatic certificates and concrete credit proofs.\newline
\href{https://github.com/VladimirReshetnikov/ProveIt/tree/56e487ce8fa4cf63ad169ffe0cd2944e9fc62dd3/Computability/BusyBeaver/Research/ZFC/k225-formal/verification/k215/simplification/proposals/automatic-certificates}{\texttt{proposals/automatic-certificates/}}.
\item ProveIt, alternative-machine index,
\href{https://github.com/VladimirReshetnikov/ProveIt/blob/56e487ce8fa4cf63ad169ffe0cd2944e9fc62dd3/Computability/BusyBeaver/Research/ZFC/candidate-k215/alternatives/README.md}{\texttt{candidate-k215/alternatives/README.md}}.
\item ProveIt, incoming quotient-detachment packet.\newline
\href{https://github.com/VladimirReshetnikov/ProveIt/blob/56e487ce8fa4cf63ad169ffe0cd2944e9fc62dd3/docs/incoming/K215_Proof_and_Quotient_Detachment_2026-10-21.zip}{Quotient-detachment research archive}.
\item ProveIt, checked framed-decoder foundation and remaining attachment
obligations.\newline
\href{https://github.com/VladimirReshetnikov/ProveIt/tree/56e487ce8fa4cf63ad169ffe0cd2944e9fc62dd3/Computability/BusyBeaver/Research/ZFC/k225-formal/verification/k215/simplification/proposals/checked-simplifications/decoder_lean}{Checked decoder library},
and the precise leaf interface in
\href{https://github.com/VladimirReshetnikov/ProveIt/blob/56e487ce8fa4cf63ad169ffe0cd2944e9fc62dd3/Computability/BusyBeaver/Research/ZFC/k225-formal/verification/k215/simplification/proposals/structural/runtime/README.md}{runtime proof factoring}.
\end{enumerate}

\end{document}
